\documentclass[12pt]{article}
\usepackage[utf8]{inputenc}
\usepackage[T1]{fontenc}
\usepackage{hyperref}
\title{Faroese Language Resources for K-12 Curriculum Development}
\author{Compiled by Claude Sonnet 5.5 (claude-sonnet-5-5) for Andr\'as Kornai}
\date{}
\begin{document}
\maketitle

\section{Introduction}
Faroese (ISO 639-3 \texttt{fao}) is an official national language of the Faroe Islands and the language of instruction in Faroese schools. It has the best-documented resource base of the four languages in this batch, though the search turned up mostly learner-oriented and reference materials rather than school textbooks.

\section{Primary Educational Resources}
\begin{itemize}
\item Faroese Online: a free introductory course with interactive exercises, and a list of recommended books and resources (written by the site's author). \url{https://faroeseonline.com/resources/}
\item Study in the Faroe Islands: LEXIN online dictionaries, designed for minority-language pupils in primary, secondary and adult education, with grammar information and speech synthesis. \url{https://studyinfaroeislands.fo/life-in-the-faroes/learn-the-language/lexin-online-encyclopedia}
\end{itemize}
No Faroese-medium school textbooks for K-12 subjects were found in the searches done.

\section{Linguistic and Reference Materials}
\begin{itemize}
\item \emph{Faroese: An Overview and Reference Grammar} (Thr\'ainsson, Petersen, Jacobsen, Svabo Hansen), an English-language reference grammar; see the listing. \url{https://nordicnames.de/wiki/Book:_Faroese_-_an_Overview_and_Reference_Grammar}
\item Freelang Faroese dictionary (Faroese--English and English--Faroese). \url{https://www.freelang.net/dictionary/faroese.php}
\end{itemize}

\section{Religious and Community Resources}
No religious or community-organisation resources were specifically identified in this search.

\section{Digital and Audio Resources}
\begin{itemize}
\item Visit Faroe Islands page on the Faroese language. \url{https://www.faroeislands.fo/arts-culture/the-faroese-language}
\item T\'orshavn municipality page on language and culture. \url{https://tvk.fo/en/living/language-culture/faroese-language/}
\end{itemize}

\section{Curriculum Development Context}
Faroese is the medium of instruction, so a curriculum would build on existing national materials. Search results point to the Sprotin dictionary publisher as the main source of authorised dictionaries, but no Sprotin URL appeared in the results, so it is not cited here. School textbook publishers and the education ministry should be consulted directly.
\end{document}